\section{데이터와 증거 범위}
\label{sec:data}

\begin{revision}

\subsection{자료 출처와 복원 범위}
\label{sec:provenance}
자연 자료는 nuScenes 기록에서 변환한 영상ㆍ자차 움직임과 사후 자동 생성한 CoC-Nusc 주석이다 \cite{caesar2020nuscenes,nvidia2026autolabeler}. 생성 과정은 기록 궤적에서 행동 종류를 계산하고 영상ㆍ행동 정보를 이용해 설명을 만든다. 따라서 CoC와 궤적의 일치는 독립적인 원인 정답이 아니다. 보관 목록에는 743장면ㆍ3,218사건ㆍ2,475인접쌍이 있으나 본문의 궤적 연결 분석은 94장면ㆍ403사건을 사용한다. 그중 CoC가 둘 이상인 90장면에서 309인접쌍을 구성했다. 사람 검토는 이 전체와 별도로 선정한 27장면 묶음이다.

CoC-Nusc의 배포 버전과 원 생성 실행별 checkpoint revision은 복원되지 않았다. 상류 README에는 모델 계열 Qwen/Qwen3.5-35B-A3B가 명시되어 있으나 개별 입력이 그 정확한 모델 revision으로 생성되었는지는 확인할 수 없다. 원 생성 prompt/config와 품질 점수 계산ㆍ상류 임계값도 미복원이다. 하류의 저장된 \texttt{quality\_pass} 불리언을 사용하는 선정 절차는 복원했다. 이번 보존 파일의 해시는 현재 입력 스냅샷을 식별하며 원 생성 이력을 새로 확정하지 않는다.

\begin{table}[t]
\centering
\color{revisionblue}
\bitablecaption{자료 출처의 확인 사실과 미복원 정보.}{Recovered provenance and information that remains unrecorded.}
\label{tab:provenance}
\footnotesize
\begin{tabularx}{\columnwidth}{L{.34\columnwidth}Y}
\toprule
항목 & 확인 범위 \\
\midrule
자연 입력 & 보관 파일ㆍ장면/사건 매핑ㆍ현재 해시 \\
생성 모델 & 상류 모델 계열 확인; 실행별 revision 미복원 \\
생성 prompt/config & 미복원; 새 prompt로 대신 채우지 않음 \\
품질 선정 & 저장된 quality\_pass 사용; 점수 계산/임계값 미복원 \\
작성 시나리오 & 원문ㆍ프레임 동시 작성, 기대 판정 실행 전 고정 \\
검사 실행 & 코드ㆍXMLㆍqueryㆍstdout/stderr 및 해시 보존 \\
\bottomrule
\end{tabularx}
\end{table}

보조 반응 분석의 선정은 $403\rightarrow290\rightarrow134$다. 품질 통과 290건 중 복합 행동 또는 방향이 불명확한 156건을 제외했다. 나머지 134건은 속도 변화 탐지 대상 130건과 사건 시점에 이미 정지ㆍ양보 조건을 만족한 4건이다. 기본 설정의 판정은 속도 반응 확인 125건, 이미 만족 4건, 제한 시간 내 미확인 5건이다. 이 선정은 309인접쌍의 문장 목록에 적용한 필터가 아니다.

\subsection{자연 검토 표본과 판정 기록}
\label{sec:selection}
후보 선택기는 이웃한 첫 문장에 \texttt{stop/yield} 계열, 둘째 문장에 \texttt{accelerate/proceed} 계열이 있는 쌍을 찾는다. 23쌍ㆍ15장면에서 장면별 가장 이른 창 하나를 선정했다. 비교군은 후보 장면을 제외하고, 첫 문장에 정지ㆍ양보 표현이 있으나 둘째 문장에는 가속ㆍ진행 표현이 없는 창이다. 장면당 가장 이른 창을 남기고, 후보와의 사건 수 차이, 정지/양보 종류 차이, 원 사건 위치 차이 순으로 대응시켰다. 동률은 seed 20260806과 창 ID의 해시로 결정했다. 비교 창을 중복 사용하지 않아 12장면이 선택되었다. 이 방식은 비무작위이며 특정 행동 전환을 더 많이 포함한다. 재분석에서 기존 27개 검토 ID와의 대응을 확인했다.

네 검토자의 원 CSV, 지침과 최종 합의 CSV를 사용했다. 지침은 시간 요구의 명시성, 요구 행동열과 문장 일관성을 구분하고, 외부 객체 상태가 필요하면 미상으로 남기도록 한다. 검토자의 전문성ㆍ교육 이력과 독립 중재자 참여는 충분한 기록이 없어 확인할 수 없다. 합의 CSV에는 27장면 모두 텍스트 모순 없음으로 저장되어 있지만 사유 메모는 0/27이다. 당시 논의 절차와 판단 변경의 원인을 재현 가능한 수준으로 복원하지 못했으며, 이번 작업에서 사람 재평가를 실시하지 않았다.

\subsection{작성 시나리오와 재현 자료}
\label{sec:artifacts}
통제 자료는 보행자, 자전거, 선행 차량, 교차 교통, 버스, 장애물의 6상황과 8변형, 중간 사건 0/2/6개의 모든 조합 144건이다. 실제 자연 문장에서 오류를 수집하거나 주입한 자료가 아니라 설계자가 영어 원문과 조건별 프레임을 함께 작성했다. 원문 구절의 시작ㆍ끝 위치, 사건ㆍ조건ㆍ의무 ID와 기대 판정이 저장되어 있다. 사례와 구현에 같은 설계자가 관여했으므로 독립 외부 평가가 아니다.

재현 자료는 프로젝트 소유 산출물 \texttt{kiee-review-revision-001} 아래에 보존했다. 공개 분류의 \texttt{controlled-scenarios-2026-10-01-001}에는 입력ㆍ최초 결과ㆍ48개 실패ㆍ모델 로그, \texttt{terminal-evidence-2026-10-01-001}에는 종료 질의 진단, \texttt{condition-model-2026-10-01-001}에는 동일 사례 재시험이 있다. 사람 원 CSVㆍ원문 및 기존 자료는 restricted 분류를 유지한다. 파일을 보존했다는 사실과 외부 공개를 완료했다는 사실은 구분한다. 이번 개정은 외부 공개를 수행하지 않았다.
\end{revision}
